\section{Computational cost and timing boundaries}
\label{app:parameter_counts}

\paragraph{Counting convention.}
Total parameters count unique trainable parameters in each analyzed
circular-cylinder specialist. Active parameters count unique trainable
parameters used by one batch-size-one prediction-path routing decision,
not their union across a batch, trajectory, or all RK4 stages. These include
the encoder and refinement blocks, executed routers, the selected group's
shared experts, channel-specific Top-2 velocity and pressure experts, and
the learned pressure correction. Fixed POD bases, projected operators,
normalization statistics, and non-trainable buffers are excluded.

\par
\begin{table}[H]
\centering
\caption{Trainable and per-decision active parameter counts for the analyzed
circular-cylinder specialists. Active counts refer to the sparse prediction
path; counts are exact integers.}
\label{tab:parameter_counts}
\begin{tabular}{lrrr}
\toprule
Specialist & Total trainable & Active prediction & Active fraction \\
\midrule
Hopf & 31,810,392 & 14,165,816 & 44.53\% \\
Periodic & 48,617,547 & 8,308,959 & 17.09\% \\
\bottomrule
\end{tabular}
\end{table}
\par

The Hopf and Periodic specialists activate 44.53\% and 17.09\% of their
trainable parameters per prediction decision, respectively. These counts
describe conditional parameter activation, not measured FLOPs, runtime,
or speedup. Complete parameter-matched counts are unavailable for the
Vanilla-FNN-MoE, DataOnly-MoE, and Global MoE ablations, so the local
architecture comparison is not interpreted as a capacity-matched study.

\paragraph{Native local latency diagnostic.}
For Circular Periodic seed 1248, batch size one and $K=48$, we use three
warm-up runs and ten timed repetitions. The native implementation
includes modal rollout, feature generation, data transfers, algebraic
pressure, all-step physical decoding, and pressure gauge correction.
It excludes physical-history encoding, E2, descriptor construction,
T2-C fusion, model/data loading, and error computation. The Galerkin
NumPy path executes on the CPU; sparse parameter activation does not
imply a proportional wall-clock saving in this implementation.

\begin{table}[H]
\centering
\caption{Local native Circular Periodic diagnostic. GPU memory is peak
allocated bytes. These times are not complete end-to-end latencies.}
\label{tab:local_latency}
\small
\begin{tabular}{@{}lrrr@{}}
\toprule
Model & Median time (s) & GPU peak (bytes) & Time / Dense\\
\midrule
Proposed & 2.1848 & 204,238,336 & 3.32\\
Dense & 0.6579 & 43,111,936 & 1.00\\
Structured & 0.8652 & 43,316,736 & 1.32\\
\bottomrule
\end{tabular}
\end{table}

Matching FOM logs alone does not supply an end-to-end speedup: the same
geometry, physical prediction interval, initialization, and compute/I/O
boundaries must also be measured for the complete ROM pipeline. No such
speedup is claimed here.

\paragraph{Paired physical-input/physical-output specialist timing.}
A separate measurement adds full-size physical-history encoding and
history right-hand-side construction to the original Periodic specialist
(seed 1248). For Circular $Re=70.314635$, the frozen first legal $K=48$
window spans $t=1384.49597$ to $1983.19690$, or 598.70093 physical time
units. Its three-frame physical input is the POD-reconstructed history
on 97,368 points; encoding is timed, but acquiring that history is not.
All 48 velocity and pressure fields are decoded and pressure is
gauge-corrected. The model and inputs are resident before timing.

On an RTX 4090 with a Xeon Platinum 8358P host, with one CPU numerical
thread, batch size one, three warm-ups, ten repetitions, and CUDA
synchronization, the median is 2.18747 s (mean 2.18729 s, sample standard
deviation 0.00462 s; range 2.18034--2.19550 s). The matched production
OpenFOAM log contains 4,692 solver steps and a ClockTime difference of
1,996 s; the largest endpoint discrepancy is $2.3\times10^{-5}$ physical
time units, attributable to stored timestamp precision. The production
solver used one process on a Core i5-14600K VMware host with 16 vCPUs and
at most five concurrent cases. Hardware identification is from the
subsequent source audit, not a contemporaneous benchmark. There is one
FOM timing realization, not ten repeated FOM runs.

The resulting ratio, $1996/2.18747\simeq912.5$, is specifically a historical
production-FOM/warm-specialist wall-time ratio. The FOM interval includes
native solver field writes but excludes meshing, initialization, periodic
detection, and subsequent VTK/NPZ conversion. ROM model/input loading,
output disk writes, E2 selection, and descriptor/T2-C operations are
excluded. Different hardware, concurrency, and I/O boundaries prevent
interpreting this ratio as a hardware-independent or complete hierarchical
pipeline speedup. A fully paired Top-2 end-to-end timing remains unmeasured.
